\section{Endpoint and extremal questions}
\label{sec:open-problems}
The preceding conclusions leave two specific questions in the obstruction
theory. They do not enter the proofs of the finite-parameter separation or
linear coercivity theorems.

\begin{problem}[The weak zero-Weyl endpoint]
Assume the original actual conformal monopole supply, $p\ge A>0$,
$a^2\ge T_+$, and, initially, $\chi(M)/A\le B$. Does there exist
$c(B)>0$ such that, for every admissible metric and every $\beta\ge2$,
\[
 \kappa_\beta:=\fint_M|\Ric+g|^2+\frac{D_\beta}{V}
 \ge c(B)\frac{(\beta-2)^2}{\beta^2}?
\]
Can the topological-density hypothesis also be removed?
\end{problem}
The $D_\beta/V$ normalization here differs from the finite-parameter
quantity $D_\beta/(\beta V)$. The unknown case is a sequence with
$\beta\to\infty$ and $\kappa_\beta\to0$ in which the exterior geometry
degenerates together with the core. A counterexample must be a sequence
of closed metrics with the full stated actual supply and budgets.
Local actual cores or formal curvature tensors alone do not supply it.
For the bounded self-dual forms associated with this limit, the Hodge
identity forces
\[
 \fint_M|\nabla\omega|^2\longrightarrow\frac83,
\]
not parallelism. The signed curvature identities in the supplements
identify an interaction that is not determined by the currently available
quadratic budgets. Their cancellation of the scalar and determinant
terms also cancels a candidate positive background term, so it cannot
be discarded in a purported contradiction.

\begin{problem}[The balanced first-order coefficient]
For the balanced family in Section~\ref{sec:variational-geometry}, let
$Q_j=A_jh_j$ and $\alpha_j=\inf_g\E(g)$. Determine
$\liminf\alpha_j/h_j$ and $\limsup\alpha_j/h_j$, with matching recovery
sequences. Equivalently, determine the first-order asymptotics of
$\inf_gP(g)$.
\end{problem}
The exact variational reduction gives
\[
 \alpha_j=4\frac{\inf_gP(g)}{8\pi^2A_j}+O(h_j^2).
\]
Consequently a coefficient of $4$ is equivalent to
$\inf_gP(g)=8\pi^2Q_j+o(Q_j)$. A lower-cost construction on the same
family is an equally legitimate alternative. The required precision is
$o(Q_j)$, not merely $o(V_j)$. A near-minimizer cannot retain a positive
density of removable over-costly cells; that fact by itself does not
exclude cheaper cells or diffuse competing geometry. The signed
characteristic formulas make this distinction explicit.

The associated existence questions are separate. At the Chern boundary
$T_+=3T_->0$, it remains to determine whether the diameter hypothesis in
the complex-hyperbolic realization theorem can be removed. On the fixed
single-Dehn-twist surface bundle $E_D$, the constructed metrics satisfy
\[
 \delta_-(G_\ell)\longrightarrow0,\qquad
 \Vol(G_\ell)=16\pi^2(p-1)(q-1)
\]
for fiber and base genera $p,q\ge2$. The value $16\pi^2$ is the
genus-two case in both factors. Whether the same smooth manifold admits
any negative Einstein metric is not decided by failure of a local repair
procedure or by nonattainment of one variational infimum.

Finally, the weaker single-solution hypothesis raises an application
question: identify a natural family not covered by the previous actual
suppliers but covered by the new hypothesis. For the stable-cohomotopy
examples, a comparison with every possible monopole class requires an
upper bound for the complete classical budget, not only the supplied
sign family. These are precise extensions of the theorems, rather than
additional assumptions in their proofs.
